\documentclass[11pt,a4paper]{article}
\usepackage[utf8]{inputenc}
\usepackage[portuguese]{babel}
\usepackage{geometry}
\usepackage{fancyhdr}
\usepackage{hyperref}
\usepackage{listings}
\usepackage{xcolor}

\geometry{left=2cm,right=2cm,top=2cm,bottom=2cm}
\setlength{\parskip}{0.4em}
\setlength{\headheight}{14pt}
\pagestyle{fancy}
\fancyhf{}
\lhead{CC8550 --- Simulacao e Teste de Software}
\rhead{Laboratorio 08}
\cfoot{\thepage}

\lstset{
    basicstyle=\ttfamily\small,
    breaklines=true,
    frame=single,
    backgroundcolor=\color{gray!10}
}

\title{
    {\Large \textbf{Centro Universitario FEI}}\\
    \vspace{0.3cm}
    {\large Ciencia da Computacao}\\[0.8cm]
    {\LARGE \textbf{Simulacao e Teste de Software}}\\[0.3cm]
    {\Large Atividade 08 --- Suite de Testes para API REST}\\[0.8cm]
    {\large \textbf{DummyJSON API}}
}
\author{
    \textbf{Aluno:} Lucas Reboucas Silva\\
    \textbf{R.A.:} 22.122.048-6\\[0.4cm]
    \textbf{Professor:} Luciano Rossi
}
\date{2026}

\begin{document}
\maketitle
\thispagestyle{empty}
\vspace{1cm}
\begin{center}
\small
Relatorio academico simplificado sobre a implementacao de testes de sistema para API REST publica.
\end{center}
\newpage

\section{Objetivo}
Construir uma suite de testes automatizados para API REST publica usando \texttt{requests}, \texttt{pytest} e \texttt{jsonschema}, validando codigos de status, schema, autenticacao, operacoes CRUD, fixture e desempenho.

\section{API escolhida e justificativa}
A API escolhida foi a \textbf{DummyJSON} (\url{https://dummyjson.com/docs}), pois disponibiliza endpoints publicos para leitura e manipulacao simulada de recursos (\texttt{/posts}) e endpoint de autenticacao (\texttt{/auth/login}), atendendo aos requisitos do laboratorio em um ambiente simples.

\section{Estrutura do projeto}
Arquivos principais do LAB-08:
\begin{itemize}
    \item \texttt{test\_api.py}: suite de testes em \texttt{pytest};
    \item \texttt{requirements.txt}: dependencias da atividade;
    \item \texttt{README.md}: instrucoes e descricao dos testes;
    \item \texttt{relatorio.tex}: documentacao deste laboratorio.
\end{itemize}

\section{Cenarios de teste implementados}
Os testes cobrem os seguintes itens:
\begin{enumerate}
    \item GET em colecao (\texttt{/posts}) com status 200 e lista nao vazia;
    \item GET em recurso existente com validacao de schema;
    \item GET em recurso inexistente (\texttt{/posts/0}) com status 404;
    \item POST para criacao (\texttt{/posts/add}) com status 201 e identificador no retorno;
    \item PATCH para atualizacao de campo (\texttt{/posts/1});
    \item DELETE (\texttt{/posts/1}) com status 200/204;
    \item envio de payload invalido com retorno 4xx;
    \item autenticacao em \texttt{/auth/login} com e sem credencial completa;
    \item uso de \texttt{@pytest.fixture} para payload reutilizavel;
    \item tempo de resposta inferior a 2 segundos.
\end{enumerate}

\section{Exemplo de comando}
\begin{lstlisting}[language=bash]
pytest test_api.py -v
\end{lstlisting}

\section{Conclusao}
A suite implementada atende os requisitos propostos da Atividade Pratica do LAB-08, validando comportamento funcional e nao funcional basico (tempo de resposta). A abordagem com API publica permitiu exercitar cenarios reais de integracao HTTP sem uso de ferramentas graficas.

\end{document}
